\section{Расчётное задание 9}

Для уравнения
\begin{equation}
	\label{eq:num_09:main}
	x^2 - 2y^2 + z^2
	+ 4xy - 8xz - 4yz
	- 14x - 4y + 14z
	+ 16
	=
	0
\end{equation}
поверхности второго порядка требуется:
\begin{enumerate}[1)]
	\item выделив квадратичную форму, найти ортогональное преобразование координат, приводящее её к каноническому виду;
	\item используя найденное преобразование и параллельный перенос, привести уравнение поверхности к каноническому виду, записать соответствующее преобразование координат, определить тип поверхности, сделать чертёж в старой системе координат.
\end{enumerate}

\begin{proofm}
	Уравнение~\eqref{eq:num_09:main} можно записать в виде
	\[
		\Vect{
			x & y & z & 1
		}
		\begin{pmatrix}
			1 & 2 & -4 & -7 \\
			2 & -2 & -2 & -2 \\
			-4 & -2 & 1 & 7  \\
			-7 & -2 & 7 & 16
		\end{pmatrix}
		\Vect{
			x \\ y \\ z \\ 1
		}
		=
		0
	\]
	или в виде
	\[
		\Vect{
			x & y & z
		}
		\begin{pmatrix}
			1 & 2 & -4  \\
			2 & -2 & -2 \\
			-4 & -2 & 1
		\end{pmatrix}
		\Vect{
			x \\ y \\ z
		}
		+
		\Vect{
			x & y & z
		}
		\Vect{
			-14 \\ -4 \\ 14
		}
		+
		16
		=
		0
		.
	\]
	Первое слагаемое --- это квадратичная форма
	\[
		q :=
		x^2 - 2y^2 + z^2
		+ 4xy - 8xz - 4yz
		.
	\]

	Приведём её к диагональному виду --- на этот раз не методом Лагранжа, а вычисляя спектр матрицы квадратичной формы~$q$. Обозначим её
	\[
		\Qe :=
		\begin{pmatrix}
			1 & 2 & -4 \\
			2 & -2 & -2 \\
			-4 & -2 & 1
		\end{pmatrix}
		.
	\]
	Найдём её спектр, решив уравнение~$\det(\Qe - tE) = 0$:
	\begin{align*}
		\det(\Qe - tE)
		&=
		\begin{determ}
			1 - t & 2 & -4 \\
			2 & -2-t & -2 \\
			-4 & -2 & 1-t
		\end{determ}
		\nlinea{=}
		(1-t)
		\begin{determ}
			-2-t & -2 \\
			-2 & 1-t
		\end{determ}
		-
		2
		\begin{determ}
			2 & -2 \\
			-4 & 1-t
		\end{determ}
		-
		4
		\begin{determ}
			2 & -2-t \\
			-4 & -2
		\end{determ}
		\nlinea{=}
		(1-t)
		\bigl(
			(-2-t)(1-t) - 4
		\bigr)
		-2
		\bigl(
			(2-2t) - 8
		\bigr)
		-
		4
		\bigl(
			-4 - 4(2+t)
		\bigr)
		\nlinea{=}
		-t^3 + 27t + 54 =
		-(t-6)(t+3)^2.
	\end{align*}
	Если~$\lambda$ --- корень характеристического полинома (который отличается только множителем~$(-1)^3 = -1$), то либо $t = \lambda_1 = 6$, либо $t = \lambda_2 = -3$ (корень кратности~$2$). Так как корней меньше размерности матрицы, то один из корней задаёт одномерное собственное подпространство, а другой --- кратный --- двумерное. Поэтому в некотором базисе~$\ss$ (мы скоро найдём, в каком именно) матрица формы~$q$ имеет вид
	\[
		\Qs =
		\left( 
			\begin{array}{c|cc}
				1 & 0 & 0 \\
				\hline
				0 & -3 & 1 \\
				0 & 0 & -3
			\end{array}
		\right)
		.
	\]
\end{proofm}
